\documentclass[11pt]{article}

% --- Folienformat 16:9 (960x540bp = 33.87x19.05cm) ---
\usepackage[paperwidth=960bp, paperheight=540bp, margin=0pt]{geometry}
\usepackage[ngerman]{babel}
\usepackage[utf8]{inputenc}
\usepackage[T1]{fontenc}
\usepackage[scaled]{helvet}
\renewcommand{\familydefault}{\sfdefault}
\usepackage{microtype}
\usepackage{xcolor}
\usepackage{tikz}
\usetikzlibrary{arrows.meta, calc, shapes.geometric, backgrounds}

\pgfdeclarelayer{bg}
\pgfdeclarelayer{edges}
\pgfdeclarelayer{labels}
\pgfsetlayers{bg,edges,main,labels}

\pagestyle{empty}
\setlength{\parindent}{0pt}
\hyphenpenalty=10000
\exhyphenpenalty=10000

% --- Reine Schwarz-Weiß / Graustufen-Palette ---
\definecolor{bgCanvas}{RGB}{255, 255, 255}
\definecolor{inkBlack}{RGB}{20, 20, 20}
\definecolor{inkMid}{RGB}{95, 95, 95}

\newcommand{\pTag}{\textbf{[+]}}
\newcommand{\mTag}{\textbf{[--]}}
\newcommand{\pmTag}{\textbf{[$\pm$]}}

\tikzset{
  every path/.style = {line cap=round, line join=round},
  % --- Knoten-Stile ---
  agoraNode/.style = {
    ellipse, draw=inkBlack, line width=1.6pt, fill=white, text=inkBlack,
    minimum width=3.20cm, minimum height=1.35cm, align=center,
    font=\Large\bfseries
  },
  coreCard/.style = {
    rounded corners=5pt, draw=inkBlack, line width=1.25pt, fill=black!3, text=inkBlack,
    minimum width=4.15cm, align=center, inner xsep=9pt, inner ysep=6pt
  },
  shellCard/.style = {
    rounded corners=4pt, draw=inkBlack!80, line width=1.05pt, fill=white, text=inkBlack,
    minimum width=3.85cm, align=center, inner xsep=8pt, inner ysep=5.5pt
  },
  % --- Einheitlicher Pfeil-Stil ---
  stdArrow/.style = {
    line width=1.05pt, draw=inkBlack!85,
    -{Stealth[length=2.3mm, width=1.55mm]},
    preaction={draw=bgCanvas, line width=3.2pt, arrows=-}
  },
  % --- Labels ---
  lbl/.style = {
    font=\fontsize{6.2}{7.1}\selectfont, text=inkBlack, fill=bgCanvas,
    inner xsep=2.4pt, inner ysep=1.2pt, align=center
  },
  lblBold/.style = {
    font=\fontsize{6.4}{7.3}\selectfont\bfseries, text=inkBlack, fill=bgCanvas,
    inner xsep=2.6pt, inner ysep=1.3pt, align=center
  },
}

\begin{document}
\begin{tikzpicture}[remember picture, overlay, shift={(current page.center)}]

  % --- Reiner weißer Hintergrund (ohne gestrichelte Ringe) ---
  \begin{pgfonlayer}{bg}
    \fill[bgCanvas] (current page.south west) rectangle (current page.north east);
  \end{pgfonlayer}

  % --- Titel (oben links) ---
  \fill[inkBlack] (-16.25, 8.95) rectangle ++(0.14, -0.92);
  \node[anchor=north west, font=\Large\bfseries, text=inkBlack, inner sep=0pt]
    at (-15.92, 8.93) {Agora: Game-Design Loops \& Schalenmodell};
  \node[anchor=north west, font=\scriptsize, text=inkMid, inner sep=0pt]
    at (-15.92, 8.30) {3 Kernelemente im Zentrum $\cdot$ 3 Schalenelemente mit gegenseitigen Abhängigkeiten};

  % ==========================================================================
  % KNOTEN
  % ==========================================================================

  % --- Zentrum: Agora ---
  \node[agoraNode] (AG) at (0.0, -0.25) {Agora};
  \draw[inkBlack!35, line width=0.65pt] (AG) ellipse (1.48cm and 0.60cm);

  % --- 3 Kernelemente ---
  \node[coreCard] (ST) at (0.0, 3.50) {
    {\normalsize\bfseries Sozialer Treffpunkt}\\[2.5pt]
    \begin{minipage}{3.2cm}
      \scriptsize
      $\bullet$ Menschen\\
      $\bullet$ Kultur\\
      $\bullet$ Freizeit
    \end{minipage}
  };

  \node[coreCard] (H) at (-5.6, -2.40) {
    {\normalsize\bfseries Handelsplatz}\\[2.5pt]
    \begin{minipage}{3.2cm}
      \scriptsize
      $\bullet$ Finanzsystem\\
      $\bullet$ Digitalität\\
      $\bullet$ Ressourcen
    \end{minipage}
  };

  \node[coreCard] (P) at (5.6, -2.40) {
    {\normalsize\bfseries Politik}\\[2.5pt]
    \begin{minipage}{3.2cm}
      \scriptsize
      $\bullet$ Exekutive\\
      $\bullet$ Legislative\\
      $\bullet$ Judikative
    \end{minipage}
  };

  % --- 3 Schalenelemente ---
  \node[shellCard] (I) at (-11.3, 3.15) {
    {\normalsize\bfseries Infrastruktur}\\[2.5pt]
    \scriptsize $\bullet$ Verkehr \enspace $\bullet$ Energie \enspace $\bullet$ Netze
  };

  \node[shellCard] (N) at (11.3, 3.15) {
    {\normalsize\bfseries Normen \& Werte}\\[2.5pt]
    \scriptsize $\bullet$ Vertrauen \enspace $\bullet$ Solidarität \enspace $\bullet$ Fairness
  };

  \node[shellCard] (U) at (0.0, -7.35) {
    {\normalsize\bfseries Unternehmen}\\[2.5pt]
    \scriptsize $\bullet$ Konzerne \enspace $\bullet$ Mittelstand \enspace $\bullet$ Start-ups
  };

  % ==========================================================================
  % KANTEN & PFEILE (einheitlicher Stil, weiche organische Kurven)
  % ==========================================================================
  \begin{pgfonlayer}{edges}

  % --- 1. Kern-Trias ---
  \draw[stdArrow] (H.north)      to[out=78, in=182, looseness=1.05] (ST.west);
  \draw[stdArrow] (ST.south west) to[out=205, in=55, looseness=1.05] (H.north east);

  \draw[stdArrow] (ST.east)      to[out=-2, in=102, looseness=1.05] (P.north);
  \draw[stdArrow] (P.north west) to[out=125, in=-25, looseness=1.05] (ST.south east);

  \draw[stdArrow] (H.east)       to[out=12, in=168, looseness=0.95] (P.west);
  \draw[stdArrow] (P.south west) to[out=-166, in=-14, looseness=0.95] (H.south east);

  % --- 2. Kern <-> Schale ---
  % Infrastruktur <-> Sozialer Treffpunkt
  \draw[stdArrow] (I.east) to[out=12, in=168] (ST.168);

  % Infrastruktur <-> Handelsplatz
  \draw[stdArrow] (I.south east) to[out=-62, in=132, looseness=1.05] (H.north west);
  \draw[stdArrow] (H.west)       to[out=162, in=-95, looseness=1.10] (I.south);

  % Sozialer Treffpunkt <-> Normen & Werte
  \draw[stdArrow] (ST.12)  to[out=12, in=168] (N.west);
  \draw[stdArrow] (N.192)  to[out=192, in=-12] (ST.350);

  % Normen & Werte <-> Politik
  \draw[stdArrow] (N.south)      to[out=-85, in=18, looseness=1.10] (P.east);
  \draw[stdArrow] (P.north east) to[out=48, in=-118, looseness=1.05] (N.south west);

  % Handelsplatz <-> Unternehmen
  \draw[stdArrow] (H.south)      to[out=-80, in=172, looseness=1.05] (U.west);
  \draw[stdArrow] (U.north west) to[out=145, in=-35, looseness=1.05] (H.south east);

  % Politik <-> Unternehmen
  \draw[stdArrow] (P.south)      to[out=-100, in=8, looseness=1.05] (U.east);
  \draw[stdArrow] (U.north east) to[out=35, in=-145, looseness=1.05] (P.south west);

  % --- 3. Schale <-> Schale (Äußerer Loop, weite elliptische Bögen) ---
  \draw[stdArrow] (N.south east) to[out=-68, in=2, looseness=1.05] (U.east);
  \draw[stdArrow] (U.west)       to[out=178, in=-112, looseness=1.05] (I.south west);
  \draw[stdArrow] (I.north east) to[out=28, in=152, looseness=0.90] (N.north west);

  \end{pgfonlayer}

  % ==========================================================================
  % LABELS
  % ==========================================================================
  \begin{pgfonlayer}{labels}
    % Kern-Trias
    \node[lblBold] at (-4.35, 1.85)  {bietet Waren, Jobs\\[-1pt]\& Güter \pTag};
    \node[lblBold] at (-2.30, 0.85)  {schafft Konsum\\[-1pt]\& Nachfrage \pTag};
    \node[lblBold] at (4.35, 1.85)   {wählt \& artikuliert\\[-1pt]Anliegen \pTag};
    \node[lblBold] at (2.30, 0.85)   {schützt Rechte\\[-1pt]\& setzt Rahmen \pmTag};
    \node[lblBold] at (0.0, -1.75)   {finanziert über Steuern \& Abgaben \pTag};
    \node[lblBold] at (0.0, -3.15)   {setzt Marktregeln \& Gesetze \pmTag};

    % Kern <-> Schale: Infrastruktur
    \node[lbl] at (-6.15, 4.15)  {schafft Zugang,\\[-1pt]Mobilität \& Raum \pTag};
    \node[lbl] at (-7.45, 0.55)  {ermöglicht Logistik\\[-1pt]\& Lieferketten \pTag};
    \node[lbl] at (-10.15, -0.55){treibt Ausbau \&\\[-1pt]Digitalisierung \pTag};

    % Kern <-> Schale: Normen & Werte
    \node[lbl] at (6.15, 4.15)   {verhandelt \& prägt\\[-1pt]Leitbilder \pTag};
    \node[lbl] at (6.15, 2.55)   {stiften Vertrauen \&\\[-1pt]Zusammenhalt \pTag};
    \node[lbl] at (10.15, -0.55) {setzen Maßstäbe\\[-1pt]für Gesetze \pTag};
    \node[lbl] at (7.45, 0.55)   {verankert Werte\\[-1pt]im Recht \pTag};

    % Kern <-> Schale: Unternehmen
    \node[lbl] at (-5.05, -5.50) {liefert Kapital,\\[-1pt]Arbeit \& Rohstoffe \pTag};
    \node[lbl] at (-2.35, -4.55) {schaffen Güter,\\[-1pt]Dienste \& Wert \pTag};
    \node[lbl] at (5.05, -5.50)  {reguliert, besteuert\\[-1pt]\& fördert \pmTag};
    \node[lbl] at (2.35, -4.55)  {betreiben Lobbyismus\\[-1pt]\& schaffen Jobs \pmTag};

    % Schale <-> Schale (Äußerer Loop)
    \node[lblBold] at (11.65, -3.55) {fordern Ethik, Fairness\\[-1pt]\& Nachhaltigkeit \mTag};
    \node[lblBold] at (-11.65, -3.55){brauchen \& verschleißen\\[-1pt]Netze/Energie \mTag};
    \node[lblBold] at (0.0, 6.45)    {vernetzt Gesellschaft \& beschleunigt Diskurs \pTag};
  \end{pgfonlayer}

  % ==========================================================================
  % LEGENDE (unten links, kompakt)
  % ==========================================================================
  \begin{scope}[shift={(-16.25, -7.45)}]
    \filldraw[fill=white, draw=inkBlack!25, rounded corners=4pt, line width=0.5pt]
      (0, 0) rectangle (5.65, -1.60);
    \node[anchor=north west, font=\scriptsize\bfseries, text=inkBlack, inner sep=0pt]
      at (0.24, -0.18) {Legende};

    \draw[inkBlack!85, line width=1.05pt, -{Stealth[length=2.1mm, width=1.45mm]}]
      (0.24, -0.68) -- (0.82, -0.68);
    \node[anchor=west, font=\tiny, text=inkBlack, inner sep=0pt]
      at (0.94, -0.68) {Systemische Wirkung / Rückkopplung};

    \node[anchor=west, font=\tiny, text=inkBlack, inner sep=0pt]
      at (0.24, -1.16) {\pTag\ verstärkend \quad \mTag\ begrenzend / Kosten \quad \pmTag\ regulierend};
  \end{scope}

\end{tikzpicture}
\end{document}
